\documentclass[UTF8,11pt]{ctexart}
\usepackage[a4paper,margin=22mm]{geometry}
\usepackage{graphicx,enumitem,amssymb,tabularx}
\usepackage[normalem]{ulem}
\usepackage{xeCJKfntef}
\setlength{\parindent}{0pt}
\setlength{\parskip}{.6em}
\sloppy
\begin{document}
\section*{2014年6月英语四级真题(第3套)}
\subsection*{2014 年6 月大学英语四级考试真题（第3 套）}


\subsection*{\textbf{Part I} \textbf{Writing} \textbf{(30 minutes)}}


Directions ：For this part, you are allowed 30 minutes to write a short essay on the following question. You should write at least \textbf{120} words but no more than \textbf{180} words.


Suppose a foreign Mend of yours is coming to visit China, what is the first place you would like to take him/her to see and why?


\subsection*{\textbf{Part II} \textbf{Listening Comprehension} \textbf{(30 minutes)}}


\subsection*{Section A}


\textbf{Directions}：In this section, you will hear 8 short conversations and 2 long conversations. At the end of each conversation\textbf{,} one or more questions will be asked about what was said. Both the conversation and the questions will be spoken only once. After each question there will be a pause. During the pause\textbf{,} you must read the four choices marked A\textbf{),}\textit{B)}, C) and D)\textbf{,} and decide which is the best answer. Then mark the corresponding letter on \textbf{Answer Sheet 1} with a single line through the centre.


1.


\begin{itemize}[leftmargin=2em,itemsep=.2em,topsep=.4em]\item[A.] He has proved to be a better reader than the woman.
\item[B.] He has difficulty understanding the book.
\item[C.] He cannot get access to the assigned book.
\item[D.] He cannot finish his assignment before the deadline.
\end{itemize}

2.


\begin{itemize}[leftmargin=2em,itemsep=.2em,topsep=.4em]\item[A.] She will drive the man to the supermarket.
\item[B.] The man should buy a car of his own.
\item[C.] The man needn’t go shopping every week.
\item[D.] She can pick the man up at the grocery store.
\end{itemize}

\noindent\begin{tabularx}{\linewidth}{@{}lXXXX@{}}3. & \textbf{A.} Get more food and drinks. & \textbf{B.} Ask his friend to come over. & \textbf{C.} Tidy up the place. & \textbf{D.} Hold a party.\\\end{tabularx}\par

4.


\begin{itemize}[leftmargin=2em,itemsep=.2em,topsep=.4em]\item[A.] The talks can be held any day except this Friday.
\item[B.] He could change his schedule to meet John Smith.
\item[C.] The first-round talks should start as soon as possible.
\item[D.] The woman should contact John Smith first.
\end{itemize}

5.


\begin{itemize}[leftmargin=2em,itemsep=.2em,topsep=.4em]\item[A.] He understands the woman’s feelings.
\item[B.] He has gone through a similar experience.
\item[C.] The woman should have gone on the field trip.
\item[D.] The teacher is just following the regulations.
\end{itemize}

6.


\begin{itemize}[leftmargin=2em,itemsep=.2em,topsep=.4em]\item[A.] She will meet the man halfway.
\item[C.] She will ask David to talk less.
\end{itemize}

\begin{itemize}[leftmargin=2em,itemsep=.2em,topsep=.4em]\item[B.] She is sorry the man will not come.
\item[D.] She has to invite David to the party.
\end{itemize}

7.


\begin{itemize}[leftmargin=2em,itemsep=.2em,topsep=.4em]\item[A.] Few students understand Prof. Johnson’s lectures.
\item[B.] Few students meet Prof. Johnson’s requirements.
\item[C.] Many students find Prof. Johnson’s lectures boring.
\end{itemize}

\begin{itemize}[leftmargin=2em,itemsep=.2em,topsep=.4em]\item[D.] Many students have dropped Prof. Johnson’s class.
\end{itemize}

\noindent\begin{tabularx}{\linewidth}{@{}lXXXX@{}}8. & \textbf{A.} Check their computer files. & \textbf{B.} Make some computations. & \textbf{C.} Study a computer program. & \textbf{D.} Assemble a computer.\\\end{tabularx}\par

\textbf{Questions 9 to 12 are based on the conversation you have just heard.}


9.


\begin{itemize}[leftmargin=2em,itemsep=.2em,topsep=.4em]\item[A.] It allows him to make a lot of friends.
\item[B.] It requires him to work long hours.
\item[C.] It enables him to apply theory to practice.
\item[D.] It helps him understand people better.
\end{itemize}

10.


\begin{itemize}[leftmargin=2em,itemsep=.2em,topsep=.4em]\item[A.] It is intellectually challenging.
\item[B.] It requires him to do washing-up all the time.
\item[C.] It exposes him to oily smoke all day long.
\item[D.] It demands physical endurance and patience.
\end{itemize}

\noindent\begin{tabularx}{\linewidth}{@{}lXXXX@{}}11. & \textbf{A.} In a hospital. & \textbf{B.} At a coffee shop. & \textbf{C.} At a laundry. & \textbf{D.} In a hotel.\\\end{tabularx}\par

12.


\begin{itemize}[leftmargin=2em,itemsep=.2em,topsep=.4em]\item[A.] Getting along well with colleagues.
\item[B.] Paying attention to every detail.
\item[C.] Planning everything in advance.
\item[D.] Knowing the needs of customers.
\end{itemize}

\textbf{Questions 13 to 15 are based on the conversation you have just heard.}


13.


\begin{itemize}[leftmargin=2em,itemsep=.2em,topsep=.4em]\item[A.] The pocket money British children get.
\item[B.] The annual inflation rate in Britain.
\item[C.] The things British children spend money on.
\item[D.] The rising cost of raising a child in Britain.
\end{itemize}

14.


\begin{itemize}[leftmargin=2em,itemsep=.2em,topsep=.4em]\item[A.] It enables children to live better.
\item[B.] It goes down during economic recession.
\item[C.] It often rises higher than inflation.
\item[D.] It has gone up 25\% in the past decade.
\end{itemize}

15.


\begin{itemize}[leftmargin=2em,itemsep=.2em,topsep=.4em]\item[A.] Save up for their future education.
\item[B.] Pay for small personal things.
\item[C.] Buy their own shoes and socks.
\item[D.] Make donations when necessary.
\end{itemize}

\subsection*{\textbf{Section B}}


Directions ：In this section\textbf{,} you will hear 3 short passages\textbf{,} at the end of each passage\textbf{,} you will hear some questions. Both the passage and questions will be spoken only once. After you hear a question, you must choose the best answer from the four choices marked A), B), C) and


\begin{itemize}[leftmargin=2em,itemsep=.2em,topsep=.4em]\item[D.] . Then mark the corresponding letter on \textbf{Answer Sheet 1} with a single line through the centre.
\end{itemize}

\subsection*{\textbf{Passage One}}


\textbf{Questions 16 to 19 are based on the passage you have just heard.}


\noindent\begin{tabularx}{\linewidth}{@{}lXXXX@{}}16. & \textbf{A.} District managers. & \textbf{B.} Regular customers. & \textbf{C.} Sales directors. & \textbf{D.} Senior clerks.\\\end{tabularx}\par

17.


\begin{itemize}[leftmargin=2em,itemsep=.2em,topsep=.4em]\item[A.] The support provided by the regular clients.
\item[B.] The initiative shown by the sales representatives.
\item[C.] The urgency of implementing the company’s plans.
\item[D.] The important part played by district managers.
\end{itemize}

18.


\begin{itemize}[leftmargin=2em,itemsep=.2em,topsep=.4em]\item[A.] Some of them were political-minded.
\item[B.] Fifty percent of them were female.
\item[C.] One third of them were senior managers.
\item[D.] Most of them were rather conservative.
\end{itemize}

19.


\begin{itemize}[leftmargin=2em,itemsep=.2em,topsep=.4em]\item[A.] He used too many quotations.
\item[B.] He was not gender sensitive.
\item[C.] He did not keep to the point.
\item[D.] He spent too much time on details.
\end{itemize}

\subsection*{\textbf{Passage Two}}


\textbf{Questions 20 to 22 are based on the passage you have just heard.}


20.


\begin{itemize}[leftmargin=2em,itemsep=.2em,topsep=.4em]\item[A.] State your problem to the head waiter.
\item[B.] Demand a discount on the dishes ordered.
\item[C.] Ask to see the manager politely but firmly.
\item[D.] Ask the name of the person waiting on you.
\end{itemize}

21.


\begin{itemize}[leftmargin=2em,itemsep=.2em,topsep=.4em]\item[A.] Your problem may not be understood correctly.
\item[B.] You don’t know if you are complaining at the right time.
\item[C.] Your complaint may not reach the person in charge.
\item[D.] You can’t tell how the person on the line is reacting.
\end{itemize}

\noindent\begin{tabularx}{\linewidth}{@{}lXXXX@{}}22. & \textbf{A.} Demand a prompt response. & \textbf{B.} Provide all the details. & \textbf{C.} Send it by express mail. & \textbf{D.} Stick to the point.\\\end{tabularx}\par

\subsection*{\textbf{Passage Three}}


\textbf{Questions 23 to 25 are based on the passage you have just heard.}


\noindent\begin{tabularx}{\linewidth}{@{}lXXXX@{}}23. & \textbf{A.} Fashion designer. & \textbf{B.} Architect. & \textbf{C.} City planner. & \textbf{D.} Engineer.\\\end{tabularx}\par

24.


\begin{itemize}[leftmargin=2em,itemsep=.2em,topsep=.4em]\item[A.] Do some volunteer work.
\item[B.] Get a well-paid part-time job.
\item[C.] Work flexible hours.
\item[D.] Go back to her previous post.
\end{itemize}

25.


\begin{itemize}[leftmargin=2em,itemsep=.2em,topsep=.4em]\item[A.] Few baby-sitters can be considered trustworthy.
\item[B.] It will add to family’s financial burden.
\item[C.] A baby-sitter is no replacement for a mother.
\item[D.] The children won’t get along with a baby-sitter.
\end{itemize}

\subsection*{\textbf{Section C}}


Directions ：In this section\textbf{,} you will hear a passage three times. When the passage is read for the first time\textbf{,} you should listen carefully for its general idea. When the passage is read for the second time\textbf{,} you are required to fill the blanks with the exact words you have just heard. Finally\textbf{,} when the passage is read for the third time\textbf{,} you should check what you have written.


Almost every child, on the first day he sets foot in a school building, is smarter, more \uline{26 }, less afraid of what he doesn’t know, better at finding and \uline{27} , more confident,\textit{ resourceful}（机敏的）,persistent and \uline{28} than he will ever be again in his schooling—or, unless he is very unusual and very lucky, for the rest of his life. Already, by paying close attention to and \uline{29} the world and people around him, and without any school-type formal instruction, he has done a task far more difficult, complicated and \uline{30} than anything he will be asked to do in school, or than any of his teachers has done for years. He has solved the \uline{31} of language. He has discovered it—babies don’t even know that language exists—and he has found out how it works and learnt to use it\uline{ 32 }. He has done it by exploring, by experimenting, by developing his own model of the grammar of language, by \uline{33} and seeing whether it works by gradually changing it and \uline{34} it until it does work. And while he has been doing this, he has been learning other things as well, including many of the \uline{35} that the schools think only they can teach him, and many that are more complicated than the ones they do try to teach him.


\subsection*{\textbf{Part}III \textbf{Reading Comprehension} \textbf{(40 minutes)}}


\subsection*{\textbf{Section A}}


Directions ：In this section\textbf{,} there is a passage with ten blanks. You are required to select one word for each blank from a list of choices given in a word bank following the passage. Read the passage through carefully before making your choices. Each choice in the bank is identified by a letter. Please mark the corresponding letter for each item on\textbf{Answer Sheet} 2 with a single line through the centre. You may not use any of the words in the bank more than once.


Questions 36 to 45 are based on the following passage.


The fact is, the world has been finding less oil than it has been using for more than twenty years now. Not only has demand been \uline{36} , but the oil we have been finding is coming from places that are \uline{37} to reach. At the same time, more of this newly \uline{38} oil is of the type that requires a greater investment to \uline{39} . And because demand for this precious resource will grow, according to some, by over 40 percent by 2025, fueling the world’s economic \uline{40} will take a lot more energy from every possible source.


The energy industry needs to get more from existing fields while continuing to search for new \uline{41} . Automakers must continue to improve fuel efficiency and perfect\textit{ hybrid} (混合动力的）vehicles. Technological improvements are needed so that wind, solar and hydrogen can be more \uline{42} parts of the energy equation. Governments need to formulate energy policies that promote \uline{43} and environmentally sound development. Consumers must be willing to pay for some of these solutions, while practicing conservation efforts of their own.


Inaction is not an \uline{44} . So let’s work together to balance this equation. We are taking some of the \uline{45} needed to get started, but we need your help to go the rest of the way.


\begin{itemize}[leftmargin=2em,itemsep=.2em,topsep=.4em]\item[A.] consequently
\item[B.] cultivate
\item[C.] declining
\item[D.] derived
\item[E.] difficult
\item[F.] discovered
\item[G.] economically
\item[H.] exception
\item[I.] feasible
\item[J.] growth
\item[K.] option
\item[L.] refine
\item[M.] reserves
\item[N.] soaring
\item[O.] steps
\end{itemize}

\subsection*{Section B}


\textbf{Directions}：In this section\textbf{,} you are going to read a passage with ten statements attached to it. Each statement contains information given in one of


the paragraphs. Identify the paragraph from which the information is derived. You may choose a paragraph more than once. Each


paragraph is marked with a letter. Answer the questions by marking the corresponding letter on \textbf{Answer Sheet 2.} I Cry, Therefore I Am


A) In 2008, at a German zoo, a\textit{ gorilla}(大猩猩）named Gana gave birth to a male infant, who died after three months. Photographs of Gana, looking stricken and inconsolable (伤心欲绝的）, attracted crowds to the zoo. Sad as the scene was, the humans, not Gana, were the only ones crying. The notion that animals can weep has no scientific basis. Years of observations by biologists Dian Fossey, who observed gorillas, and Jane Goodall, who worked with\textit{ chimpanzees}(黑猩猩), could not prove that animals cry tears from emotion.


B) It’s true that many animals shed tears, especially in response to pain. Tears protect the eye by keeping it moist. But crying as an expression of feeling is unique to humans and has played an essential role in human evolution and the development of human cultures.


C) Within two days an infant can imitate sad and happy faces. If an infant does not cry out, it is unlikely to get the attention it needs to survive. Around 3-4 months, the relationship between the human infant and its environment takes on a more organized communicative role, and tearful crying begins to serve interpersonal purposes: the search for comfort and \textit{pacification} (抚慰). As we get older, crying becomes a tool of social interaction: grief and joy, shame and pride, fear and manipulation.


D) Tears are as universal as laughter, and grief is more complex than joy. But although we all cry, we do so in different ways. Women cry more frequently and intensely than men, especially when exposed to emotional events. Like crying, depression is, around the world, more commonly seen in women than in men. One explanation might be that women, who despite decades of social advances still suffer from economic inequality,\textit{ discrimination} (歧视) and even violence, might have more to cry about. Men not only cry for shorter periods than women, but they also are less inclined to explain their tears, usually shed them more quietly, and tend more frequently to apologize when they cry openly. Men, like women, report crying at the death of a loved one and in response to a moving religious experience. They are more likely than women to cry when their core identities—as providers and protectors, as fathers and fighters—are questioned.


E) People who score on personality tests as more sympathetic cry more than those who are more rigid or have more self-control. Frequency of crying varies widely: some shed tears at any novel or movie, others only a handful of times in their lives. Crying in response to stress and conflict in the home, or after emotional\textit{ trauma} (创伤), lasts much longer than tears induced by everyday sadness—which in turn last longer than tears of delight and joy.


F) Sadness is our primary association with crying, but the fact is that people report feeling happier after crying. Surveys estimate that 85\% of women and 73\% of men report feeling better after shedding tears. Surprisingly, crying is more commonly associated with minor forms of depression than with major depression involving suicidal thoughts.


G) People widely report that crying relieves tension, restores emotional balance and provides “catharsis,” a washing out of bad feelings. The term “catharsis” has religious implications of removing evil and sin; it’s no surprise that religious ceremonies are, around the world, one of the main settings for the release of tears.


H) Crying is a nearly universal sign of grief, though some mourners report that, despite genuine sorrow, they cannot shed tears—sometimes even for years after their loved one has gone. Unlike today, when the privacy of grief is more respected, the public or ceremonial shedding of tears, at the graveside of a spouse or the funeral of a king or queen, was once considered socially or even politically essential.


I) Crying has also served other social purposes. Rousseau wrote in his\textit{ Confessions} that while he considered tears the most powerful expression of love, he also just liked to cry over nothing.


J) The association of tears with art has ancient roots. The classic Greek tragedies of the fifth century


B)


C) were primarily


celebrations of gods. Tragedies, like poetry and music, were staged religious events. Even then it was recognized that crying in response to drama brought pleasure.


K) I have argued that there are\textit{ neurobiological}(神经生物方面的）associations linking the arts and mood disorders. When I lecture on crying, I ask my audience to let me know, by a show of hands, which art forms most move them to tears. About 80\% say music, followed closely by novels (74\% ), but then the figures fall sharply, to 43\% , for poetry, and 10 -22\% for paintings, sculpture and architecture.


L) The physical act of crying is mainly one of breathing in air, which is why we choke up when we weep. This suggests to language scientists that emotional crying evolved before language, perhaps explaining why tears communicate states of mind and feelings that are often so difficult to express in words. Of course, from an evolutionary perspective, recognition of emotion (usually through facial gesture) was essential for survival.


M) The earliest humans arrived several million years ago, but only 150,000 to 200,000 years ago, did cultures, language, religion and the arts arise. Along the way, tears became more than a biological necessity to\textit{ lubricate}（润滑）the eye and developed into a sign of intense emotion and a signal of social bonding. The development of self-consciousness and the notion of individual identity, or ego; storytelling about the origins of the world, the creation of humanity and life after death; and the ability to feel others, sadness—all were critical parts of the neurobiological changes that made us human.


N) More recently, we’ve learned from neuroscience that certain brain\textit{ circuits}（回路）are\textit{ activated}(激活), rapidly and unconsciously, when we see another in emotional distress. In short, our brain evolved circuits to allow us to experience sympathy, which in turn made civilization, and an ethics based on sympathy, possible. So the next time you reach a tissue box, or sob on a Mend’s shoulder, or shed tears at the movies, stop and reflect on why we cry and what it means to cry. Because ultimately, while we love to cry, we also cry to love.


46. Nowadays people respect the privacy of grief more than in the past.


47. Infants cry to attract attention for survival.


48. There is no scientific evidence as yet that animals can shed tears from emotion.


49. Tears can perform certain communicative functions which words cannot.


50. Our ability to experience sympathy is essential to the development of civilization.


51. People are more inclined to cry when suffering minor forms of depression.


52. Sometimes people cannot cry despite genuine grief.


53. In humans, long history, tears have developed an essential role in social relationships.


54. Men are less likely to give reasons for their tears.


55. Crying has long been associated with art.


\subsection*{Section C}


\textbf{Directions}：There are 2 passages in this section. Each passage is followed by some questions or unfinished statements. For each of them there are


four choices marked A), B), C) and D). You should decide on the best choice and mark the corresponding letter on\textbf{ Answer}


\textbf{Sheet 2} with a single line through the centre.


\subsection*{Passage One}


\textbf{Questions 56 to 60 are based on the following passage.}


Hospitals, hoping to curb medical error, have invested heavily to put computers, smartphones and other devices into the hands of medical staff for instant access to patient data, drug information and case studies.


But like many cures, this solution has come with an unintended side effect: doctors and nurses can be focused on the screen and not the patient, even during moments of critical care. A poll showed that half of medical technicians had admitted texting during a procedure.


This phenomenon has set off an intensifying discussion at hospitals and medical schools about a problem perhaps best described as “distracted doctoring.” In response, some hospitals have begun limiting the use of electronic devices in critical settings, while schools have started reminding medical students to focus on patients instead of devices.


“You justify carrying devices around the hospital to do medical records, but you can surf the Internet or do Facebook, and sometimes Facebook is more tempting,” said Dr. Peter Papadakos at the University of Rochester Medical Center.


“My\textit{ gut feeling}（本能的感觉）is lives are in danger,” said Dr. Papadakos. “We’re not educating people about the problem, and ifs getting worse.”


A survey of 439 medical technicians found that 55 percent of technicians who monitor bypass machines acknowledged that they had talked on cellphones during heart surgery. Half said they had texted while in surgery. The study concluded, “ Such distractions have the potential to be disastrous.”


Medical professionals have always faced interruptions from cellphones, and multitasking is simply a fact of life for many medical jobs. What has changed, say doctors, especially younger ones, is that they face increasing pressure to interact with their devices.


The pressure stems from a\textit{ mantra}(信条）of modem medicine that patient care must be “data driven,” and informed by the latest, instantly accessible information. By many accounts, the technology has helped reduce medical error by providing instant access to patient data or prescription details.


Dr. Peter Carmel, president of the American Medical Association, said technology “offers great potential in health c a r e , ” b u t he added that doctors, first priority should be with the patient.


56. Why do hospitals equip their staff with computers, smartphones and other devices?


\begin{itemize}[leftmargin=2em,itemsep=.2em,topsep=.4em]\item[A.] To reduce medical error.
\item[B.] To cope with emergencies.
\item[C.] To facilitate administration.
\item[D.] To simplify medical procedures.
\end{itemize}

57. What does the author refer to by “distracted doctoring”?


\begin{itemize}[leftmargin=2em,itemsep=.2em,topsep=.4em]\item[A.] The disservice done by modem devices to doctors, nurses, as well as patients.
\item[B.] The tendency of medical institutions encouraging the use of modem devices.
\item[C.] The problem of devices preventing doctors from focusing on their patients.
\item[D.] The phenomenon of medical staff attending to personal affairs while working.
\end{itemize}

58. What does Dr. Peter Papadakos worry about?


\begin{itemize}[leftmargin=2em,itemsep=.2em,topsep=.4em]\item[A.] Medical students are not adequately trained to use modem technology.
\item[B.] Doctors’ interaction with their devices may endanger patients, lives.
\item[C.] Doctors are relying too heavily on modem electronic technology.
\item[D.] Pressures on the medical profession may become overwhelming.
\end{itemize}

59. Why do doctors feel increasing pressure to use modem devices?


\begin{itemize}[leftmargin=2em,itemsep=.2em,topsep=.4em]\item[A.] Patients trust doctors who use modem technology.
\item[B.] Use of modem devices adds hospitals’ revenues.
\item[C.] Data is given too much importance in patient care.
\item[D.] Patients’ data has to be revised from time to time.
\end{itemize}

60. What is Peter Carmel’s advice to doctors?


\begin{itemize}[leftmargin=2em,itemsep=.2em,topsep=.4em]\item[A.] They follow closely the advances in medical science.
\item[B.] They focus their attention on the patient’s condition.
\item[C.] They observe hospital rules and regulations.
\item[D.] They make the best use of modem devices.
\end{itemize}

\subsection*{Passage Two}


Questions 61 to 65 are based on the following passage.


I have closely watched my generation, known as The Millennials, for 29 years now. Joel Stein wrote an extensive piece on


Millennials and he remains rather optimistic about our potential.


I hesitate to share his optimism because of\textit{ paradox} (矛盾的现象）we seem to exhibit, namely, that there are more avenues for us to entertain ourselves than ever before, yet we are more bored than ever before.


Entertainment has never been more varied. We have more cable channels, television shows, and movies than ever before. Internet providers allow instant viewing of almost any movie or television program ever created. Social drinking and partying are also widely available for Millennials. Every generation develops these habits at a certain age, but Millennials seem to be extending this phase of life as they postpone marriage.


Some of this is undoubtedly due to The Great Recession. Milleimials are having a difficult time finding jobs; only 47 percent of 16-to-24-year-olds are employed, the smallest share since government started recording data in 1948.


But do Millennials respond to these economic troubles by doing whatever it takes to make ends meet? Hardly. In fact, of the four generations Pew Research has data for, the Milennial generation does not cite\textit{ work} \textit{ethic}（勤奋工作）as distinctive of itself. Millennials want to save the world, but they sit and wait for that world-changing opportunity to be handed to them. Instead of working 2-3 jobs, launching a business, or doing what it takes to succeed, they retreat. Millennials may be the first generation to have a lower standard of living than their parents, but with this response to\textit{ adversity}(逆境), perhaps deservingly so.


Much ink has been spilled in management books discussing how to get the most out of these youths in the workplace. Largely, they come to the same conclusion: Millennials are entitled, over-confident, and expect too much too quickly. We should not be surprised. Today’s young adults were raised by parents who made sure to boost their self-esteem at every turn, telling them they could achieve whatever they set their minds to, and handing out prizes for the sixth place.


61. What does the author of the passage think of Millennials?


\begin{itemize}[leftmargin=2em,itemsep=.2em,topsep=.4em]\item[A.] They show little interest in entertainment.
\item[B.] They are not confident about their ability.
\item[C.] They enjoy an easy life due to high technology.
\item[D.] They may not have bright prospects for success.
\end{itemize}

62. How do Millennials feel about their life?


\begin{itemize}[leftmargin=2em,itemsep=.2em,topsep=.4em]\item[A.] They can hardly do anything about it.
\item[B.] There is little in it to get excited about.
\item[C.] It is not as good as their parents’.
\item[D.] It is full of opportunities for success.
\end{itemize}

63. In what way are Millennials different from previous generations according to Pew Research?


\begin{itemize}[leftmargin=2em,itemsep=.2em,topsep=.4em]\item[A.] They spend less time socializing.
\item[B.] They are indifferent to others.-
\item[C.] They do not value hard work.
\item[D.] They are more independent.
\end{itemize}

64. What should Millennials do according to the author?


\begin{itemize}[leftmargin=2em,itemsep=.2em,topsep=.4em]\item[A.] Remain optimistic in face of adversity.
\item[B.] Start a business as early as possible.
\item[C.] Make full use of new opportunities.
\item[D.] Take action to change their situation.
\end{itemize}

65. Why are Millennials over-confident about themselves?


\begin{itemize}[leftmargin=2em,itemsep=.2em,topsep=.4em]\item[A.] They have been spoiled by their parents.
\item[B.] They can always get whatever they expect.
\item[C.] They are misguided by management books.
\item[D.] They think they are young and energetic.
\end{itemize}

\subsection*{\textbf{Part IV} \textbf{Translation} \textbf{(30 minutes)}}


\textbf{Directions}：For this part\textbf{,} you are allowed 30 minutes to translate a passage from Chinese into English. You should write your answer on\textbf{Answer} \textbf{Sheet} 2.


为了促进教育公平，中国已投入360 亿元,用于改善农村地区教育设施和加强中西部地区农村\CJKunderline{义务教育} (compulsory education)。这些资金用于改善教学设施、购买书籍，使16 万多所中小学受益。资金还用于购置音乐和绘画器材。现在农村和山区的儿童可以与沿海城市的儿童一样上音乐和绘画课。一些为接受更好教育而转往城市上学的学生如今又回到了本地农村学校就读。

\end{document}